% ============================================================================
% SECTION 1.5: DEMOCRATIZATION THROUGH USE CASES (REVISED - Dual Barriers)
% ============================================================================

\section{Democratization Through Adaptive Routing}
\label{sec:use_cases}

To ground our technical contributions in real-world impact, we illustrate how \ours{} removes two barriers that prevent access to AI capabilities: \emph{cost barriers} (frontier models are prohibitively expensive) and \emph{expertise barriers} (existing routing tools require ML knowledge and extensive setup). Each scenario demonstrates how BanditGPT makes adaptive routing accessible to users who previously lacked either budget or technical infrastructure.

\subsection{Student Projects: Learning Without Prerequisites}

\paragraph{The Dual Barrier.} A student building a paper summarization tool (5,000 abstracts) faces prohibitive costs (\$21.90 for GPT-4o-only) and lacks the hundreds to thousands of labeled examples typically required to calibrate FrugalGPT~\cite{chen2023frugalgpt}. The student has neither budget nor training data.

\paragraph{With BanditGPT.} The student installs the library via pip, downloads pre-trained priors (1MB), and writes 10 lines of Python:

\begin{verbatim}
from banditgpt import Router
router = Router(mode="standard", lambda_cost=10)
results = [router.query(abstract) for abstract in dataset]
\end{verbatim}

The router autonomously identifies that 72\% of abstracts can be handled by Gemini~Flash~Lite (\$0.10/1k), routing only complex technical content to GPT-4o. Total cost: \textbf{\$3.50} (84\% reduction). Setup time: \textbf{5 minutes}. Required expertise: \textbf{None}---no labeled data, no scorer design, no chain configuration.

\paragraph{Accessibility Impact.} By removing both economic and expertise barriers, \ours{} expands access to AI education from elite institutions (which provide API budgets and graduate-level ML courses) to community colleges and online learners who lack both resources. Students learn to build AI-augmented applications through direct experimentation rather than through prerequisite coursework.

\subsection{Independent Researchers: Methods Without Infrastructure}

\paragraph{The Expertise Barrier Amplified.} A social science researcher analyzing 10,000 interview transcripts for thematic patterns faces not only cost barriers (\$438 for GPT-4o-only) but also methodological barriers. Implementing FrugalGPT requires:
\begin{itemize}[leftmargin=*,nosep]
    \item Annotating 1,000+ examples with "correct" thematic codes (weeks of work)
    \item Training a domain-specific scoring function (requires ML expertise)
    \item Selecting and configuring model chains (requires understanding of LLM capabilities)
    \item Re-running calibration if coding schemes evolve (expensive iteration)
\end{itemize}

For researchers without computational backgrounds, these requirements are insurmountable regardless of budget.

\paragraph{With BanditGPT.} The researcher provides a simple prompt template and cost constraint:

\begin{verbatim}
router = Router(mode="hybrid", lambda_cost=5)
codes = router.query_batch(transcripts, 
                           template="Identify themes: {text}")
\end{verbatim}

The system starts with pre-trained priors (no user calibration), autonomously discovers that 88\% of transcripts involve straightforward thematic tagging (routed to Nova-Lite), and reserves GPT-4o for nuanced, ambiguous passages. Total cost: \textbf{\$14.20} (68\% reduction). Setup complexity: \textbf{Zero ML expertise required}.

\paragraph{Methodological Shift.} This enables \emph{exploratory analysis} previously reserved for computational social science labs with dedicated ML engineers. Researchers iterate on prompt designs, test multiple coding schemes, and refine their approach---all within discretionary budgets and without needing to learn contextual bandit theory or train BERT-based regressors.

\subsection{Startups: Deployment Without Specialists}

\paragraph{The Infrastructure Barrier.} An early-stage legal tech startup (3 engineers, no ML team) wants to add AI-powered contract review. Beyond the \$52k annual inference cost (GPT-4o-only), implementing FrugalGPT requires:
\begin{itemize}[leftmargin=*,nosep]
    \item Hiring an ML engineer to design and train scoring functions (\$150k+ salary)
    \item Collecting and annotating calibration datasets for each contract type
    \item Maintaining a separate ML pipeline for model updates and chain reconfigurations
    \item Monitoring scorer drift as legal language evolves
\end{itemize}

For a pre-revenue startup with 18-month runway, this infrastructure cost (\$225k total: \$150k salary + \$75k operational overhead) exceeds the inference cost savings, making adaptive routing economically infeasible despite technical merit.

\paragraph{With BanditGPT.} The startup deploys without dedicated ML infrastructure:

\begin{verbatim}
router = Router(mode="hybrid", lambda_cost=3)
# Automatic routing, no scorer training required
analysis = router.query(contract, task="clause_review")
\end{verbatim}

Standard mode reduces inference costs to \textbf{\$8.4k annually} (84\% reduction). Hybrid mode provides high-assurance routing (\$16k annually) for risk-sensitive legal queries. Crucially, deployment requires \textbf{zero ML expertise}---the founding team implements adaptive routing with their existing engineering skillset, eliminating the need to hire specialists.

\paragraph{Product Viability.} By removing the infrastructure barrier, \ours{} converts AI from "requires dedicated ML team" to "add via standard engineering workflow." This enables small teams to compete with incumbents who have ML departments, democratizing access to AI-augmented product features.

\subsection{Enterprises: Scale Without Experts}

\paragraph{The Organizational Barrier.} A Fortune 500 customer support organization considering LLM-powered automation (10M queries/year) faces organizational friction. While a centralized ML team \emph{could} implement FrugalGPT, the workflow requires:
\begin{itemize}[leftmargin=*,nosep]
    \item ML team collects calibration data from support team (cross-org coordination)
    \item Weeks of iteration to design domain-specific scoring functions
    \item Ongoing maintenance as support categories evolve (dependencies on ML team)
    \item Separate deployments for each product line (no transfer learning)
\end{itemize}

This organizational overhead delays deployment by 6--12 months, during which executives lose patience and the initiative stalls.

\paragraph{With BanditGPT.} The support engineering team deploys directly without ML dependencies:

\begin{verbatim}
router = Router(mode="hybrid", 
                lambda_cost=1,  # Enterprise: prioritize reliability
                lambda_latency=0.5)
\end{verbatim}

Standard mode reduces costs from \$43.8M to \textbf{\$7.0M annually} (84\% reduction). Hybrid mode (\$13.4M) provides selective verification for high-stakes queries (billing disputes, account closures). Deployment time: \textbf{2 weeks} (vs.\ 6--12 months for FrugalGPT setup). Required cross-org dependencies: \textbf{None}---support engineers own the entire pipeline.

\paragraph{Organizational Transformation.} By removing ML team dependencies, \ours{} enables decentralized AI adoption. Support, sales, HR, and legal departments deploy independently, accelerating time-to-value from quarters to weeks.

\subsection{Cross-Cutting Themes: Two Barriers, One Solution}

Three patterns emerge across these use cases:

\paragraph{1. Ease of Use Unlocks Adoption.}
Cost reduction alone is insufficient if setup requires ML expertise. Students lack training data. Researchers lack computational backgrounds. Startups lack ML teams. By providing \emph{immediate usability} through pre-trained priors and autonomous learning, we expand the user base from "ML practitioners" to "anyone who can write Python."

\paragraph{2. Iteration Speed Enables Creativity.}
When setup takes days and requires specialists, users cannot afford experimentation. With 5-minute deployment and zero calibration requirements, users shift from "Will this work?" to "What else can I try?" This cognitive shift drives innovation in application design.

\paragraph{3. Quality Preservation Builds Trust.}
Users hesitate to adopt cost optimization out of fear that quality will degrade. Our hybrid mode (98\% reliability, \S\ref{sec:rq3}) provides a tunable dial: students maximize volume in standard mode; enterprises ensure reliability in hybrid mode. This flexibility accommodates diverse risk profiles without requiring users to choose between "cheap" and "good."

\subsection{Comparison: Setup Requirements}

\begin{table}[t]
\centering
\small
\caption{\textbf{Accessibility Comparison.} We prioritize ease of deployment to expand the user base beyond ML experts.}
\label{tab:accessibility_comparison}
\resizebox{\columnwidth}{!}{%
\begin{tabular}{lccc}
\toprule
\textbf{Requirement} & \textbf{FrugalGPT} & \textbf{RouteLLM} & \textbf{BanditGPT} \\
\midrule
Calibration Data & 500--2k examples & 1k--5k pairs & 0 (shippable priors) \\
Annotated Labels & Yes (ground truth) & Yes (preferences) & No \\
Scorer Training & Yes (BERT finetune) & Yes (classifier) & No \\
Model Chain Design & Manual (expert knowledge) & N/A (2 models) & Autonomous \\
Setup Time & Days & Hours & Minutes \\
ML Expertise Required & High & Medium & None \\
\midrule
Target User & ML teams & ML practitioners & Anyone \\
\bottomrule
\end{tabular}%
}
\end{table}

Table~\ref{tab:accessibility_comparison} summarizes the operational requirements. By eliminating calibration data, scorer training, and manual chain design, \ours{} reduces the barrier from "requires ML team" to "requires pip install." This shift expands adaptive routing from specialized deployments to mainstream adoption.

\paragraph{Design Philosophy.} These scenarios informed \ours{}'s architecture: pre-trained shippable priors eliminate calibration requirements; autonomous online learning removes expertise dependencies; tunable $\lambda$ parameters support diverse operational contexts; open-source release ensures zero licensing barriers. The technical contributions in \S\ref{sec:method} exist not as algorithmic novelty for its own sake, but as mechanisms for democratizing access to adaptive routing.

